% chapter07.tex — Three Hard Theorems
% Source: sources/chapters/ch07-three-hard-theorems.pdf (book pp. 122–132).
% Chapter-PDF leaves 127/128 and 131/132 are swapped; transcribed in book order.

\spivakchapter{7}{Three Hard Theorems}

This chapter is devoted to three theorems about continuous functions, and some
of their consequences. The proofs of the three theorems themselves will not be
given until the next chapter, for reasons which are explained at the end of this
chapter.

\begin{theorem}
If $f$ is continuous on $[a,b]$ and $f(a)<0<f(b)$, then there is some $x$ in
$[a,b]$ such that $f(x)=0$.
\end{theorem}

(Geometrically, this means that the graph of a continuous function which starts
below the horizontal axis and ends above it must cross this axis at some point,
as in Figure~1.)
\begin{center}
\figtex{assets/geometric/ch07-fig01}\\[0.6em]
\figlabel{1}
\end{center}

\begin{theorem}
If $f$ is continuous on $[a,b]$, then $f$ is bounded above on $[a,b]$, that is,
there is some number $N$ such that $f(x)\leq N$ for all $x$ in $[a,b]$.
\end{theorem}

(Geometrically, this theorem means that the graph of $f$ lies below some line
parallel to the horizontal axis, as in Figure~2.)
\begin{center}
\figtex{assets/geometric/ch07-fig02}\\[0.6em]
\figlabel{2}
\end{center}

\begin{theorem}
If $f$ is continuous on $[a,b]$, then there is some number $y$ in $[a,b]$ such
that $f(y)\geq f(x)$ for all $x$ in $[a,b]$ (Figure~3).
\end{theorem}
\begin{center}
\figtex{assets/geometric/ch07-fig03}\\[0.6em]
\figlabel{3}
\end{center}

These three theorems differ markedly from the theorems of Chapter~6. The
hypotheses of those theorems always involved continuity at a single point, while
the hypotheses of the present theorems require continuity on a whole interval
$[a,b]$---if continuity fails to hold at a single point, the conclusions may fail.
For example, let $f$ be the function shown in Figure~4,
\[
f(x)=\begin{cases}
-1, & 0\leq x<\sqrt{2}\\
1,  & \sqrt{2}\leq x\leq 2.
\end{cases}
\]
\begin{center}
\figtex{assets/geometric/ch07-fig04}\\[0.6em]
\figlabel{4}
\end{center}
Then $f$ is continuous at every point of $[0,2]$ except $\sqrt{2}$, and
$f(0)<0<f(2)$, but there is no point $x$ in $[0,2]$ such that $f(x)=0$; the
discontinuity at the single point $\sqrt{2}$ is sufficient to destroy the
conclusion of Theorem~1.

Similarly, suppose that $f$ is the function shown in Figure~5,
\[
f(x)=\begin{cases}
1/x, & x\neq 0\\
0,   & x=0.
\end{cases}
\]
\begin{center}
\figtex{assets/plots/ch07-fig05}\\[0.6em]
\figlabel{5}
\end{center}
Then $f$ is continuous at every point of $[0,1]$ except $0$, but $f$ is not
bounded above on $[0,1]$. In fact, for any number $N>0$ we have
$f(1/2N)=2N>N$.

This example also shows that the closed interval $[a,b]$ in Theorem~2 cannot be
replaced by the open interval $(a,b)$, for the function $f$ is continuous on
$(0,1)$, but is not bounded there.

Finally, consider the function shown in Figure~6,
\[
f(x)=\begin{cases}
x^2, & x<1\\
0,   & x\geq 1.
\end{cases}
\]
\begin{center}
\figtex{assets/geometric/ch07-fig06}\\[0.6em]
\figlabel{6}
\end{center}
On the interval $[0,1]$ the function $f$ is bounded above, so $f$ does satisfy
the conclusion of Theorem~2, even though $f$ is not continuous on $[0,1]$. But
$f$ does not satisfy the conclusion of Theorem~3---there is no $y$ in $[0,1]$
such that $f(y)\geq f(x)$ for all $x$ in $[0,1]$; in fact, it is certainly not
true that $f(1)\geq f(x)$ for all $x$ in $[0,1]$ so we cannot choose $y=1$, nor
can we choose $0\leq y<1$ because $f(y)<f(x)$ if $x$ is any number with
$y<x<1$.

This example shows that Theorem~3 is considerably stronger than Theorem~2.
Theorem~3 is often paraphrased by saying that a continuous function on a closed
interval ``takes on its maximum value'' on that interval.

As a compensation for the stringency of the hypotheses of our three theorems,
the conclusions are of a totally different order than those of previous
theorems. They describe the behavior of a function, not just near a point, but
on a whole interval; such ``global'' properties of a function are always
significantly more difficult to prove than ``local'' properties, and are
correspondingly of much greater power. To illustrate the usefulness of
Theorems~1, 2, and~3, we will soon deduce some important consequences, but it
will help to first mention some simple generalizations of these theorems.

\begin{theorem}
If $f$ is continuous on $[a,b]$ and $f(a)<c<f(b)$, then there is some $x$ in
$[a,b]$ such that $f(x)=c$.
\end{theorem}

\begin{proof}
Let $g=f-c$. Then $g$ is continuous, and $g(a)<0<g(b)$. By Theorem~1, there is
some $x$ in $[a,b]$ such that $g(x)=0$. But this means that $f(x)=c$.
\end{proof}

\begin{theorem}
If $f$ is continuous on $[a,b]$ and $f(a)>c>f(b)$, then there is some $x$ in
$[a,b]$ such that $f(x)=c$.
\end{theorem}

\begin{proof}
The function $-f$ is continuous on $[a,b]$ and $-f(a)<-c<-f(b)$. By Theorem~4
there is some $x$ in $[a,b]$ such that $-f(x)=-c$, which means that $f(x)=c$.
\end{proof}

Theorems~4 and~5 together show that $f$ takes on any value between $f(a)$ and
$f(b)$. We can do even better than this: if $c$ and $d$ are in $[a,b]$, then
$f$ takes on any value between $f(c)$ and $f(d)$. The proof is simple: if, for
example, $c<d$, then just apply Theorems~4 and~5 to the interval $[c,d]$.
Summarizing, if a continuous function on an interval takes on two values, it
takes on every value in between; this slight generalization of Theorem~1 is
often called the Intermediate Value Theorem.

\begin{theorem}
If $f$ is continuous on $[a,b]$, then $f$ is bounded below on $[a,b]$, that is,
there is some number $N$ such that $f(x)\geq N$ for all $x$ in $[a,b]$.
\end{theorem}

\begin{proof}
The function $-f$ is continuous on $[a,b]$, so by Theorem~2 there is a number
$M$ such that $-f(x)\leq M$ for all $x$ in $[a,b]$. But this means that
$f(x)\geq -M$ for all $x$ in $[a,b]$, so we can let $N=-M$.
\end{proof}

Theorems~2 and~6 together show that a continuous function $f$ on $[a,b]$ is
bounded on $[a,b]$, that is, there is a number $N$ such that $\abs{f(x)}\leq N$
for all $x$ in $[a,b]$. In fact, since Theorem~2 ensures the existence of a
number $N_1$ such that $f(x)\leq N_1$ for all $x$ in $[a,b]$, and Theorem~6
ensures the existence of a number $N_2$ such that $f(x)\geq N_2$ for all $x$ in
$[a,b]$, we can take $N=\max(\abs{N_1},\abs{N_2})$.

\begin{theorem}
If $f$ is continuous on $[a,b]$, then there is some $y$ in $[a,b]$ such that
$f(y)\leq f(x)$ for all $x$ in $[a,b]$.
(A continuous function on a closed interval takes on its minimum value on that
interval.)
\end{theorem}

\begin{proof}
The function $-f$ is continuous on $[a,b]$; by Theorem~3 there is some $y$ in
$[a,b]$ such that $-f(y)\geq -f(x)$ for all $x$ in $[a,b]$, which means that
$f(y)\leq f(x)$ for all $x$ in $[a,b]$.
\end{proof}

Now that we have derived the trivial consequences of Theorems~1, 2, and~3, we
can begin proving a few interesting things.

\begin{theorem}
Every positive number has a square root. In other words, if $\alpha>0$, then
there is some number $x$ such that $x^2=\alpha$.
\end{theorem}

\begin{proof}
Consider the function $f(x)=x^2$, which is certainly continuous. Notice that the
statement of the theorem can be expressed in terms of $f$: ``the number $\alpha$
has a square root'' means that $f$ takes on the value $\alpha$. The proof of
this fact about $f$ will be an easy consequence of Theorem~4.

There is obviously a number $b>0$ such that $f(b)>\alpha$ (as illustrated in
Figure~7); in fact, if $\alpha>1$ we can take $b=\alpha$, while if $\alpha<1$ we
can take $b=1$. Since $f(0)<\alpha<f(b)$, Theorem~4 applied to $[0,b]$ implies
that for some $x$ (in $[0,b]$), we have $f(x)=\alpha$, i.e., $x^2=\alpha$.
\end{proof}
\begin{center}
\figtex{assets/plots/ch07-fig07}\\[0.6em]
\figlabel{7}
\end{center}

Precisely the same argument can be used to prove that a positive number has an
$n$th root, for any natural number $n$. If $n$ happens to be odd, one can do
better: \emph{every} number has an $n$th root. To prove this we just note that
if the positive number $\alpha$ has the $n$th root $x$, i.e., if $x^n=\alpha$,
then $(-x)^n=-\alpha$ (since $n$ is odd), so $-\alpha$ has the $n$th root $-x$.
The assertion, that for odd $n$ any number $\alpha$ has an $n$th root, is
equivalent to the statement that the equation
\[
x^n-\alpha=0
\]
has a root if $n$ is odd. Expressed in this way the result is susceptible of
great generalization.

\begin{theorem}
If $n$ is odd, then any equation
\[
x^n + a_{n-1}x^{n-1} + \cdots + a_0 = 0
\]
has a root.
\end{theorem}

\begin{proof}
We obviously want to consider the function
\[
f(x)=x^n + a_{n-1}x^{n-1} + \cdots + a_0;
\]
we would like to prove that $f$ is sometimes positive and sometimes negative.
The intuitive idea is that for large $\abs{x}$, the function is very much like
$g(x)=x^n$ and, since $n$ is odd, this function is positive for large positive
$x$ and negative for large negative $x$. A little algebra is all we need to
make this intuitive idea work.

The proper analysis of the function $f$ depends on writing
\[
f(x)=x^n + a_{n-1}x^{n-1} + \cdots + a_0
= x^n\Bigl(1 + \frac{a_{n-1}}{x} + \cdots + \frac{a_0}{x^n}\Bigr).
\]
Note that
\[
\Biggl\lvert \frac{a_{n-1}}{x} + \frac{a_{n-2}}{x^2} + \cdots + \frac{a_0}{x^n}
\Biggr\rvert
\leq \frac{\abs{a_{n-1}}}{\abs{x}} + \cdots + \frac{\abs{a_0}}{\abs{x^n}}.
\]
Consequently, if we choose $x$ satisfying
\begin{equation*}
\abs{x}>1,\ 2n\abs{a_{n-1}},\ \dots,\ 2n\abs{a_0}, \tag{$*$}
\end{equation*}
then $\abs{x^k}>\abs{x}$ and
\[
\frac{\abs{a_{n-k}}}{\abs{x^k}}
< \frac{\abs{a_{n-k}}}{\abs{x}}
< \frac{\abs{a_{n-k}}}{2n\abs{a_{n-k}}}
= \frac{1}{2n},
\]
so
\[
\Biggl\lvert \frac{a_{n-1}}{x} + \frac{a_{n-2}}{x^2} + \cdots + \frac{a_0}{x^n}
\Biggr\rvert
\leq \underbrace{\frac{1}{2n}+\cdots+\frac{1}{2n}}_{n\text{ terms}}
= \frac12.
\]
In other words,
\[
-\frac12 \leq \frac{a_{n-1}}{x} + \cdots + \frac{a_0}{x^n} \leq \frac12,
\]
which implies that
\[
\frac12 \leq 1 + \frac{a_{n-1}}{x} + \cdots + \frac{a_0}{x^n}.
\]
Therefore, if we choose an $x_1>0$ which satisfies ($*$), then
\[
\frac{(x_1)^n}{2}
\leq (x_1)^n\Biggl(1 + \frac{a_{n-1}}{x_1} + \cdots + \frac{a_0}{(x_1)^n}\Biggr)
= f(x_1),
\]
so that $f(x_1)>0$. On the other hand, if $x_2<0$ satisfies ($*$), then
$(x_2)^n<0$ and
\[
\frac{(x_2)^n}{2}
\geq (x_2)^n\Biggl(1 + \frac{a_{n-1}}{x_2} + \cdots + \frac{a_0}{(x_2)^n}\Biggr)
= f(x_2),
\]
so that $f(x_2)<0$.

Now applying Theorem~1 to the interval $[x_2,x_1]$ we conclude that there is an
$x$ in $[x_2,x_1]$ such that $f(x)=0$.
\end{proof}

Theorem~9 disposes of the problem of odd degree equations so happily that it
would be frustrating to leave the problem of even degree equations completely
undiscussed. At first sight, however, the problem seems insuperable. Some
equations, like $x^2-1=0$, have a solution, and some, like $x^2+1=0$, do
not---what more is there to say? If we are willing to consider a more general
question, however, something interesting \emph{can} be said. Instead of trying
to solve the equation
\[
x^n + a_{n-1}x^{n-1} + \cdots + a_0 = 0,
\]
let us ask about the possibility of solving the equations
\[
x^n + a_{n-1}x^{n-1} + \cdots + a_0 = c
\]
for all possible numbers $c$. This amounts to allowing the constant term $a_0$
to vary. The information which can be given concerning the solution of these
equations depends on a fact which is illustrated in Figure~8.
\begin{center}
\figtex{assets/geometric/ch07-fig08}\\[0.6em]
\figlabel{8}
\end{center}
The graph of the function $f(x)=x^n+a_{n-1}x^{n-1}+\cdots+a_0$, with $n$ even,
contains, at least the way we have drawn it, a lowest point. In other words,
there is a number $y$ such that $f(y)\leq f(x)$ for all numbers $x$---the
function $f$ takes on a minimum value, not just on each closed interval, but on
the whole line. (Notice that this is false if $n$ is odd.) The proof depends on
Theorem~7, but a tricky application will be required. We can apply Theorem~7 to
any interval $[a,b]$, and obtain a point $y_0$ such that $f(y_0)$ is the
minimum value of $f$ on $[a,b]$; but if $[a,b]$ happens to be the interval
shown in Figure~8, for example, then the point $y_0$ will not be the place
where $f$ has its minimum value for the whole line. In the next theorem the
entire point of the proof is to choose an interval $[a,b]$ in such a way that
this cannot happen.

\begin{theorem}
If $n$ is even and $f(x)=x^n+a_{n-1}x^{n-1}+\cdots+a_0$, then there is a number
$y$ such that $f(y)\leq f(x)$ for all $x$.
\end{theorem}

\begin{proof}
As in the proof of Theorem~9, if
\[
M=\max\bigl(1,\ 2n\abs{a_{n-1}},\ \dots,\ 2n\abs{a_0}\bigr),
\]
then for all $x$ with $\abs{x}\geq M$, we have
\[
\frac12 \leq 1 + \frac{a_{n-1}}{x} + \cdots + \frac{a_0}{x^n}.
\]
Since $n$ is even, $x^n\geq 0$ for all $x$, so
\[
\frac{x^n}{2}
\leq x^n\Biggl(1 + \frac{a_{n-1}}{x} + \cdots + \frac{a_0}{x^n}\Biggr)
= f(x),
\]
provided that $\abs{x}\geq M$. Now consider the number $f(0)$. Let $b>0$ be a
number such that $b^n\geq 2f(0)$ and also $b>M$. Then, if $x\geq b$, we have
(Figure~9)
\[
f(x)\geq \frac{x^n}{2}\geq \frac{b^n}{2}\geq f(0).
\]
\begin{center}
\figtex{assets/geometric/ch07-fig09}\\[0.6em]
\figlabel{9}
\end{center}
Similarly, if $x\leq -b$, then
\[
f(x)\geq \frac{x^n}{2}\geq \frac{(-b)^n}{2}=\frac{b^n}{2}\geq f(0).
\]
Summarizing:
\[
\text{if } x\geq b \text{ or } x\leq -b, \text{ then } f(x)\geq f(0).
\]
Now apply Theorem~7 to the function $f$ on the interval $[-b,b]$. We conclude
that there is a number $y$ such that
\begin{align*}
&(1) & &\text{if } -b\leq x\leq b, \text{ then } f(y)\leq f(x).
\end{align*}
In particular, $f(y)\leq f(0)$. Thus
\begin{align*}
&(2) & &\text{if } x\leq -b \text{ or } x\geq b, \text{ then }
  f(x)\geq f(0)\geq f(y).
\end{align*}
Combining (1) and (2) we see that $f(y)\leq f(x)$ for all $x$.
\end{proof}

Theorem~10 now allows us to prove the following result.

\begin{theorem}
Consider the equation
\begin{equation*}
x^n + a_{n-1}x^{n-1} + \cdots + a_0 = c, \tag{$*$}
\end{equation*}
and suppose $n$ is even. Then there is a number $m$ such that ($*$) has a
solution for $c\geq m$ and has no solution for $c<m$.
\end{theorem}

\begin{proof}
Let $f(x)=x^n+a_{n-1}x^{n-1}+\cdots+a_0$ (Figure~10).
\begin{center}
\figtex{assets/geometric/ch07-fig10}\\[0.6em]
\figlabel{10}
\end{center}
According to Theorem~10 there is a number $y$ such that $f(y)\leq f(x)$ for all
$x$. Let $m=f(y)$. If $c<m$, then the equation ($*$) obviously has no
solution, since the left side always has a value $\geq m$. If $c=m$, then
($*$) has $y$ as a solution. Finally, suppose $c>m$. Let $b$ be a number such
that $b>y$ and $f(b)>c$. Then $f(y)=m<c<f(b)$. Consequently, by Theorem~4,
there is some number $x$ in $[y,b]$ such that $f(x)=c$, so $x$ is a solution of
($*$).
\end{proof}

These consequences of Theorems~1, 2, and~3 are the only ones we will derive now
(these theorems will play a fundamental role in everything we do later,
however). Only one task remains---to prove Theorems~1, 2, and~3.
Unfortunately, we cannot hope to do this---on the basis of our present
knowledge about the real numbers (namely, P1--P12) a proof is
\emph{impossible}. There are several ways of convincing ourselves that this
gloomy conclusion is actually the case. For example, the proof of Theorem~8
relies only on the proof of Theorem~1; if we could prove Theorem~1, then the
proof of Theorem~8 would be complete, and we would have a proof that every
positive number has a square root. As pointed out in Part~I, it is impossible
to prove this on the basis of P1--P12. Again, suppose we consider the function
\[
f(x)=\frac{1}{x^2-2}.
\]
If there were no number $x$ with $x^2=2$, then $f$ would be continuous, since
the denominator would never $=0$. But $f$ is not bounded on $[0,2]$. So
Theorem~2 depends essentially on the existence of numbers other than rational
numbers, and therefore on some property of the real numbers other than
P1--P12.

Despite our inability to prove Theorems~1, 2, and~3, they are certainly results
which we want to be true. If the pictures we have been drawing have any
connection with the mathematics we are doing, if our notion of continuous
function corresponds to any degree with our intuitive notion, Theorems~1, 2,
and~3 have got to be true. Since a proof of any of these theorems must require
some new property of $\mathbf{R}$ which has so far been overlooked, our present
difficulties suggest a way to discover that property: let us try to construct a
proof of Theorem~1, for example, and see what goes wrong.

One idea which seems promising is to locate the first point where $f(x)=0$,
that is, the smallest $x$ in $[a,b]$ such that $f(x)=0$. To find this point,
first consider the set $A$ which contains all numbers $x$ in $[a,b]$ such that
$f$ is negative on $[a,x]$. In Figure~11, $x$ is such a point, while $x'$ is
not. The set $A$ itself is indicated by a heavy line. Since $f$ is negative at
$a$, and positive at $b$, the set $A$ contains some points greater than $a$,
while all points sufficiently close to $b$ are not in $A$. (We are here using
the continuity of $f$ on $[a,b]$, as well as Problem~6-16.)
\begin{center}
\figtex{assets/geometric/ch07-fig11}\\[0.6em]
\figlabel{11}
\end{center}

Now suppose $\alpha$ is the smallest number which is greater than all members
of $A$; clearly $a<\alpha<b$. We claim that $f(\alpha)=0$, and to prove this we
only have to eliminate the possibilities $f(\alpha)<0$ and $f(\alpha)>0$.

Suppose first that $f(\alpha)<0$. Then, by Theorem~6-3, $f(x)$ would be less
than $0$ for all $x$ in a small interval containing $\alpha$, in particular for
some numbers bigger than $\alpha$ (Figure~12); but this contradicts the fact
that $\alpha$ is bigger than every member of $A$, since the larger numbers
would also be in $A$. Consequently, $f(\alpha)<0$ is false.
\begin{center}
\figtex{assets/geometric/ch07-fig12}\\[0.6em]
\figlabel{12}
\end{center}

On the other hand, suppose $f(\alpha)>0$. Again applying Theorem~6-3, we see
that $f(x)$ would be positive for all $x$ in a small interval containing
$\alpha$, in particular for some numbers smaller than $\alpha$ (Figure~13).
This means that these smaller numbers are all \emph{not} in $A$. Consequently,
one could have chosen an even smaller $\alpha$ which would be greater than all
members of $A$. Once again we have a contradiction; $f(\alpha)>0$ is also
false. Hence $f(\alpha)=0$ and, we are tempted to say, Q.E.D.
\begin{center}
\figtex{assets/geometric/ch07-fig13}\\[0.6em]
\figlabel{13}
\end{center}

We know, however, that something must be wrong, since no new properties of
$\mathbf{R}$ were ever used, and it does not require much scrutiny to find the
dubious point. It is clear that we can choose a number $\alpha$ which is
greater than all members of $A$ (for example, we can choose $\alpha=b$), but it
is not so clear that we can choose a \emph{smallest} one. In fact, suppose $A$
consists of all numbers $x\geq 0$ such that $x^2<2$. If the number $\sqrt{2}$
did not exist, there would not be a least number greater than all the members
of $A$; for any $y>\sqrt{2}$ we chose, we could always choose a still smaller
one.

Now that we have discovered the fallacy, it is almost obvious what additional
property of the real numbers we need. All we must do is say it properly and use
it. That is the business of the next chapter.

\subsection*{PROBLEMS}

\pnum{1.} For each of the following functions, decide which are bounded above
or below on the indicated interval, and which take on their maximum or minimum
value. (Notice that $f$ \emph{might} have these properties even if $f$ is not
continuous, and even if the interval is not a closed interval.)
\begin{romans}
\item $f(x)=x^2$ on $(-1,1)$.
\item $f(x)=x^3$ on $(-1,1)$.
\item $f(x)=x^2$ on $\mathbf{R}$.
\item $f(x)=x^2$ on $[0,\infty)$.
\item $\displaystyle f(x)=\begin{cases}
  x^2, & x\leq a\\
  a+2, & x>a
\end{cases}$
on $(-a-1,a+1)$. (We assume $a>-1$, so that $-a-1<a+1$; it will be necessary
to consider several possibilities for $a$.)
\item $\displaystyle f(x)=\begin{cases}
  x^2, & x<a\\
  a+2, & x\geq a
\end{cases}$
on $[-a-1,a+1]$. (Again assume $a>-1$.)
\item $\displaystyle f(x)=\begin{cases}
  0, & x\text{ irrational}\\
  1/q, & x=p/q\text{ in lowest terms}
\end{cases}$
on $[0,1]$.
\item $\displaystyle f(x)=\begin{cases}
  1, & x\text{ irrational}\\
  1/q, & x=p/q\text{ in lowest terms}
\end{cases}$
on $[0,1]$.
\item $\displaystyle f(x)=\begin{cases}
  1, & x\text{ irrational}\\
  -1/q, & x=p/q\text{ in lowest terms}
\end{cases}$
on $[0,1]$.
\item $\displaystyle f(x)=\begin{cases}
  x, & x\text{ rational}\\
  0, & x\text{ irrational}
\end{cases}$
on $[0,a]$.
\item $f(x)=\sin^2\bigl(\cos x+\sqrt{a+a^2}\bigr)$ on $[0,a^3]$.
\item $f(x)=[x]$ on $[0,a]$.
\end{romans}

\pnum{2.} For each of the following polynomial functions $f$, find an integer
$n$ such that $f(x)=0$ for some $x$ between $n$ and $n+1$.
\begin{romans}
\item $f(x)=x^3-x+3$.
\item $f(x)=x^5+5x^4+2x+1$.
\item $f(x)=x^5+x+1$.
\item $f(x)=4x^2-4x+1$.
\end{romans}

\pnum{3.} Prove that there is some number $x$ such that
\begin{romans}
\item $\displaystyle x^{179}+\frac{163}{1+x^2+\sin^2 x}=119$.
\item $\sin x=x-1$.
\end{romans}

\pnum{4.} This problem is a continuation of Problem~3-7.
\begin{letters}
\item If $n-k$ is even, and $\geq 0$, find a polynomial function of degree $n$
with exactly $k$ roots.
\item A root $a$ of the polynomial function $f$ is said to have
\textbf{multiplicity $m$} if $f(x)=(x-a)^m g(x)$, where $g$ is a polynomial
function that does \emph{not} have $a$ as a root. Let $f$ be a polynomial
function of degree $n$. Suppose that $f$ has $k$ roots, counting
multiplicities, i.e., suppose that $k$ is the sum of the multiplicities of all
the roots. Show that $n-k$ is even.
\end{letters}

\pnum{5.} Suppose that $f$ is continuous on $[a,b]$ and that $f(x)$ is always
rational. What can be said about $f$?

\pnum{6.} Suppose that $f$ is a \emph{continuous} function on $[-1,1]$ such that
$x^2+(f(x))^2=1$ for all $x$. (This means that $(x,f(x))$ always lies on the
unit circle.) Show that either $f(x)=\sqrt{1-x^2}$ for all $x$, or else
$f(x)=-\sqrt{1-x^2}$ for all $x$.

\pnum{7.} How many continuous functions $f$ are there which satisfy
$(f(x))^2=x^2$ for all $x$?

\pnum{8.} Suppose that $f$ and $g$ are continuous, that $f^2=g^2$, and that
$f(x)\neq 0$ for all $x$. Prove that either $f(x)=g(x)$ for all $x$, or else
$f(x)=-g(x)$ for all $x$.

\pnum{9.}
\begin{letters}
\item Suppose that $f$ is continuous, that $f(x)=0$ only for $x=a$, and that
$f(x)>0$ for some $x>a$ as well as for some $x<a$. What can be said about
$f(x)$ for all $x\neq a$?
\item Again assume that $f$ is continuous and that $f(x)=0$ only for $x=a$, but
suppose, instead, that $f(x)>0$ for some $x>a$ and $f(x)<0$ for some $x<a$. Now
what can be said about $f(x)$ for $x\neq a$?
\item[\textbf{*(c)}] Discuss the sign of $x^3+x^2 y+xy^2+y^3$ when $x$ and $y$
are not both $0$.
\end{letters}

\pnum{10.} Suppose $f$ and $g$ are continuous on $[a,b]$ and that $f(a)<g(a)$,
but $f(b)>g(b)$. Prove that $f(x)=g(x)$ for some $x$ in $[a,b]$. (If your proof
isn't very short, it's not the right one.)

\pnum{11.} Suppose that $f$ is a continuous function on $[0,1]$ and that
$f(x)$ is in $[0,1]$ for each $x$ (draw a picture). Prove that $f(x)=x$ for
some number $x$.

\pnum{12.}
\begin{letters}
\item Problem~11 shows that $f$ intersects the diagonal of the square in
Figure~14 (solid line). Show that $f$ must also intersect the other (dashed)
diagonal.
\begin{center}
\figtex{assets/geometric/ch07-fig14}\\[0.6em]
\figlabel{14}
\end{center}
\item Prove the following more general fact: If $g$ is continuous on $[0,1]$
and $g(0)=0$, $g(1)=1$ or $g(0)=1$, $g(1)=0$, then $f(x)=g(x)$ for some $x$.
\end{letters}

\pnum{13.}
\begin{letters}
\item Let $f(x)=\sin 1/x$ for $x\neq 0$ and let $f(0)=0$. Is $f$ continuous on
$[-1,1]$? Show that $f$ satisfies the conclusion of the Intermediate Value
Theorem on $[-1,1]$; in other words, if $f$ takes on two values somewhere on
$[-1,1]$, it also takes on every value in between.
\item[\textbf{*(b)}] Suppose that $f$ satisfies the conclusion of the
Intermediate Value Theorem, and that $f$ takes on each value \emph{only once}.
Prove that $f$ is continuous.
\item[\textbf{*(c)}] Generalize to the case where $f$ takes on each value only
finitely many times.
\end{letters}

\pnum{14.} If $f$ is a continuous function on $[0,1]$, let $\lVert f\rVert$ be
the maximum value of $\abs{f}$ on $[0,1]$.
\begin{letters}
\item Prove that for any number $c$ we have $\lVert cf\rVert=\abs{c}\cdot\lVert f\rVert$.
\item[\textbf{*(b)}] Prove that $\lVert f+g\rVert\leq\lVert f\rVert+\lVert g\rVert$.
Give an example where $\lVert f+g\rVert\neq\lVert f\rVert+\lVert g\rVert$.
\item Prove that $\lVert h-f\rVert\leq\lVert h-g\rVert+\lVert g-f\rVert$.
\end{letters}

\pnum{*15.} Suppose that $\phi$ is continuous and
\[
\lim_{x\to\infty}\phi(x)/x^n=0=\lim_{x\to-\infty}\phi(x)/x^n.
\]
\begin{letters}
\item Prove that if $n$ is odd, then there is a number $x$ such that
$x^n+\phi(x)=0$.
\item Prove that if $n$ is even, then there is a number $y$ such that
$y^n+\phi(y)\leq x^n+\phi(x)$ for all $x$.
\end{letters}
Hint: Of which proofs does this problem test your understanding?

\pnum{*16.}
\begin{letters}
\item Suppose that $f$ is continuous on $(a,b)$ and
$\lim_{x\to a^+}f(x)=\lim_{x\to b^-}f(x)=\infty$. Prove that $f$ has a minimum
on all of $(a,b)$.
\item Prove the corresponding result when $a=-\infty$ and/or $b=\infty$.
\end{letters}

\pnum{*17.} Let $f$ be any polynomial function. Prove that there is some number
$y$ such that $\abs{f(y)}\leq\abs{f(x)}$ for all $x$.

\pnum{*18.} Suppose that $f$ is a continuous function with $f(x)>0$ for all
$x$, and $\lim_{x\to\infty}f(x)=0=\lim_{x\to-\infty}f(x)$. (Draw a picture.)
Prove that there is some number $y$ such that $f(y)\geq f(x)$ for all $x$.

\pnum{*19.}
\begin{letters}
\item Suppose that $f$ is continuous on $[a,b]$, and let $x$ be any number.
Prove that there is a point on the graph of $f$ which is closest to $(x,0)$; in
other words there is some $y$ in $[a,b]$ such that the distance from $(x,0)$ to
$(y,f(y))$ is $\leq$ distance from $(x,0)$ to $(z,f(z))$ for all $z$ in
$[a,b]$. (See Figure~15.)
\begin{center}
\figtex{assets/geometric/ch07-fig15}\\[0.6em]
\figlabel{15}
\end{center}
\item Show that this same assertion is not necessarily true if $[a,b]$ is
replaced by $(a,b)$ throughout.
\item Show that the assertion \emph{is} true if $[a,b]$ is replaced by
$\mathbf{R}$ throughout.
\item In cases (a) and (c), let $g(x)$ be the minimum distance from $(x,0)$ to
a point on the graph of $f$. Prove that $g(y)\leq g(x)+\abs{x-y}$, and conclude
that $g$ is continuous.
\item Prove that there are numbers $x_0$ and $x_1$ in $[a,b]$ such that the
distance from $(x_0,0)$ to $(x_1,f(x_1))$ is $\leq$ the distance from
$(x_0',0)$ to $(x_1',f(x_1'))$ for any $x_0'$, $x_1'$ in $[a,b]$.
\end{letters}

\pnum{20.}
\begin{letters}
\item Suppose that $f$ is continuous on $[0,1]$ and $f(0)=f(1)$. Let $n$ be any
natural number. Prove that there is some number $x$ such that
$f(x)=f(x+1/n)$, as shown in Figure~16 for $n=4$. Hint: Consider the function
$g(x)=f(x)-f(x+1/n)$; what would be true if $g(x)\neq 0$ for all $x$?
\begin{center}
\figtex{assets/geometric/ch07-fig16}\\[0.6em]
\figlabel{16}
\end{center}
\item[\textbf{*(b)}] Suppose $0<a<1$, but that $a$ is not equal to $1/n$ for
any natural number $n$. Find a function $f$ which is continuous on $[0,1]$ and
which satisfies $f(0)=f(1)$, but which does not satisfy $f(x)=f(x+a)$ for any
$x$.
\end{letters}

\pnum{*21.}
\begin{letters}
\item Prove that there does not exist a continuous function $f$ defined on
$\mathbf{R}$ which takes on every value exactly twice. Hint: If $f(a)=f(b)$ for
$a<b$, then either $f(x)>f(a)$ for all $x$ in $(a,b)$ or $f(x)<f(a)$ for all
$x$ in $(a,b)$. Why? In the first case all values close to $f(a)$, but slightly
larger than $f(a)$, are taken on somewhere in $(a,b)$; this implies that
$f(x)<f(a)$ for $x<a$ and $x>b$.
\item Refine part (a) by proving that there is no continuous function $f$ which
takes on each value either $0$ times or $2$ times, i.e., which takes on exactly
twice each value that it does take on. Hint: The previous hint implies that $f$
has either a maximum or a minimum value (which must be taken on twice). What
can be said about values close to the maximum value?
\item Find a continuous function $f$ which takes on every value exactly $3$
times. More generally, find one which takes on every value exactly $n$ times,
if $n$ is odd.
\item Prove that if $n$ is even, then there is no continuous $f$ which takes on
every value exactly $n$ times. Hint: To treat the case $n=4$, for example, let
$f(x_1)=f(x_2)=f(x_3)=f(x_4)$. Then either $f(x)>0$ for all $x$ in two of the
three intervals $(x_1,x_2)$, $(x_2,x_3)$, $(x_3,x_4)$, or else $f(x)<0$ for all
$x$ in two of these three intervals.
\end{letters}
